\section{Introduction}
\label{sec:intro}

As large language models (LLMs) take on engineering problems, it becomes essential to evaluate how
they reason, not only what they answer.
Engineering reasoning requires physically grounded, multi-step problem solving, in which scientific
principles, quantitative modeling, and practical constraints must converge, and in which
intermediate steps are as consequential as the final result.
A correct final answer can still follow flawed reasoning~\citep{cobbe2021gsm8k}, so scoring the
answer alone is not enough.

\begin{figure}[t]
    \centering
    \includegraphics[width=\linewidth]{figs/engtrace-example-template.pdf}
    \caption{A simplified \ourdataset template (chemical branch), the instance it generates, and its
        gold reasoning trace, computed together with the answer. We check model responses against
        this trace.}
    \label{fig:example}
\end{figure}

Current benchmarks rarely evaluate this process.
General benchmarks such as \textsc{MMLU}, \textsc{MATH}, and \textsc{HumanEval} test knowledge,
mathematical problem solving, or code generation in
isolation~\citep{hendrycks2021mmlu, hendrycks2021math, chen2021humaneval}, and most engineering
benchmarks grade the final answer rather than the steps that produce
it~\citep{du2025supergpqa, zhou2025engibench}.
Those that do assess intermediate steps mostly rely on LLM judges or human graders, or cover a
single domain~\citep{syed2024transportbench, zheng2026tpscalcbench, duzkar2026thermoqa}.
To address this gap, we introduce \ourdataset, a symbolic benchmark for verifiable process
supervision of engineering reasoning.
Following symbolic-template benchmarks in mathematics and
finance~\citep{mirzadeh2024gsmsymbolic, xie2025finchain}, \ourdataset represents each problem as a
template: expert-certified code that samples physically grounded parameters, computes the answer,
and emits a gold reasoning trace from the same computation.
\autoref{fig:example} shows such a template, an instance it generates, and the gold trace; because
the trace comes from the computation itself, every intermediate quantity of the derivation is known
exactly.

These gold traces make reasoning checkable by computation.
An expert-validated evaluator scores each response deterministically wherever a check exists: it
verifies final answers of six kinds, measures the coverage of the gold trace's intermediate
quantities, which we call \emph{milestones}, and recomputes the arithmetic that the response
displays.
For the residual milestones that deterministic matching cannot find, the evaluator consults an LLM
judge from outside the evaluated model families, because LLM judges can favor their own outputs and
those of related models~\citep{panickssery2024selfrecognition, spiliopoulou2025playfavorites}.
Three experts from the template's own branch certify each template, and we draw the evaluated
instances with a private seed, so that no instance was public when the models ran.
\ourdataset is a depth-oriented evaluation of a curriculum-based subset of engineering: its problems
are examination-style by construction, as verifiable gold traces require, which leaves ill-posed
problems and real-world artifacts out of scope.

In summary, our contributions are as follows:
\begin{itemize}
    \item We introduce \ourdataset, 150 expert-certified symbolic templates spanning five
          engineering branches, 15 domains, and 42 areas, whose code-computed gold traces make every
          intermediate quantity of an engineering derivation checkable.
    \item We design and validate an evaluator that checks reasoning by computation wherever
          possible: against step-level labels from 15 engineering experts on responses of five
          models outside the evaluation, its final-answer check agrees on over 90\% of responses and
          its milestone coverage reaches an F1 of 0.96, and we identify the errors that no
          deterministic check can detect.
    \item We evaluate eleven open-weights and closed LLMs on 2,250 unseen instances and find that
          the final answer alone does not tell the strongest models apart: final-answer accuracy
          ranges from 0.81 to 0.98, with the five strongest within 0.01 of one another, while
          milestone coverage ranks the models differently. Calculation slips dominate errors on easy
          problems, while wrong formulas or principles appear mostly on intermediate and advanced
          problems.
\end{itemize}
